\subsection{The complete observed search and conditional expected bands}
\label{sec:v5-band-results}

The full search runs from 19 to 250~MeV. At each mass it includes every campaign
whose declared search interval covers that point: 2015 alone at 19--38~MeV,
2015+2016 at 39--49~MeV, all three at 50--100~MeV, 2016+2021 at 101--180~MeV,
and 2021 alone at 181--250~MeV. This membership rule is fixed by the search
ranges; it never selects the dataset giving the smallest observed limit.
Figure~\ref{fig:v5-union-bands} places the nominal local asymptotic $p_0$ directly
below the corresponding observed limit and conditional expected bands, followed
by an aligned significance panel. Black gives the local reference and red its
Sidak correction for this complete search.

The bands use 300 background-only draws per mass from the
mass-local GP prediction. Each draw first samples a correlated latent
background and then Poisson counts; constituent draws are shared among the
combinations at that mass. The GP state is held fixed. Different masses do not
form one coherent full-spectrum experiment. The resulting quantiles describe
pointwise variation under that conditional model; they do not calibrate the
\CLs{} sampling distributions or supply a global significance.

\fig{0.98}{../figures/v5_total_limits_with_local_pvalues.pdf}{Complete 19--250~MeV
observed search, with the preassigned active datasets shown above the curves.
The upper panel shows 90\% \CLs{} limits and the median, central 68\% and central
95\% ranges of the 300 conditional background-only toys per mass, with newly computed
bands at the added 2015 coordinates.
The next two panels show the fixed-mass asymptotic $p_0$ and
its normal-quantile significance, with the local reference in black and the
Sidak reference in red. The latter uses $N_{\rm eff}=35.381$ for the
complete search; its value at 66~MeV is $Z_{\rm Sidak}=1.305$,
compared with local $Z=2.765$. The observed
polyline connects every adjacent tested mass, including membership transitions;
no numerical endpoint is changed by those connections. The physical
visible-branching correction is applied once above the dimuon threshold.
These are pointwise bands; 2016-inclusive values retain the numerical exception.}
{fig:v5-union-bands}

\paragraph{Reading the significance panels.}
We use the one-sided normal convention $Z=\Phi^{-1}(1-p)$ and display
$Z^+=\max(0,Z)$, with integer guides on the vertical axis. Thus a probability
of at least one half appears at zero; this convention avoids interpreting a
negative normal quantile as positive discovery evidence. It does not change
the probability. The signed fit coordinate $r$, which can describe a deficit,
is shown separately in Section~\ref{sec:v5-global-results}.

The Sidak curves apply Eq.~\eqref{eq:sidak-global} to each local threshold
using one fixed $N_{\rm eff}$ for the entire corresponding search.
The inherited resolution-spacing prescription sums
$\Delta m/(2.25\,\bar\sigma_m)$ between adjacent mass hypotheses, uses the
narrowest active campaign resolution, and caps the result between one and
the number of tested masses. It gives 14.521, 16.062 and 26.402 for the
individual 2015, 2016 and 2021 searches. This is an analytic reference for
the look-elsewhere penalty, not a calibration from complete simulated scans.
Table~\ref{tab:v502-sidak-lookup} shows how the same local significance maps
to different corrections as the search range and resolution change.

\begin{table}[htbp]\centering\small
\begin{tabular}{rrrrrr}\toprule
Local $Z$ & Local $p$ & 2015 & 2016 & 2021 & Complete search\\
 & & \multicolumn{4}{c}{Sidak $Z^+$}\\\midrule
1 & 0.159 & 0.000 & 0.000 & 0.000 & 0.000\\
2 & 0.0228 & 0.571 & 0.499 & 0.112 & 0.000\\
3 & 0.00135 & 2.066 & 2.024 & 1.811 & 1.678\\
4 & 3.17e-05 & 3.314 & 3.286 & 3.143 & 3.056\\
5 & 2.87e-07 & 4.457 & 4.435 & 4.327 & 4.262\\
\bottomrule\end{tabular}

\caption{Local-to-Sidak significance conversion for the declared individual
and complete searches. These are illustrative thresholds passed through the
fixed analytic formula, not additional observed peaks. The final four columns
use the $Z^+$ display convention; signed quantiles and probabilities are
retained in the accompanying numerical tables.}
\label{tab:v502-sidak-lookup}
\end{table}

\fig{0.98}{../figures/v5_observed_coupling_overlay.pdf}{Observed coupling limits
from the three individual released samples and the full maximal-available-dataset
combination of Fig.~\ref{fig:v5-union-bands}. Only observed limits are shown.
The combination is obtained from a shared-coupling likelihood with the fixed
membership rule; it is not the envelope of the individual curves.}
{fig:v5-observed-overlay}

\FloatBarrier

\paragraph{The 2015 90--100 MeV extension.}
All 2015-containing curves now include the ten added integer hypotheses.
The 14--135~MeV training support is retained and the 90~MeV kernel parameters
are held fixed above that endpoint. The template width follows the existing
linear response, rising from 4.70 to 5.23~MeV across the added interval.
The shaded range has new observed limits, 300 conditional toys per mass and
scope, and the saved 95~MeV recovery check. Appendix~\ref{app:v504-extension}
records the calculation and its validation limits. Historical displays are
retained in Appendix~\ref{app:v504-historical-displays}.
